\documentclass[10pt,a4paper]{amsart}
\usepackage{amsmath,amssymb,amsthm,graphicx}
\usepackage{verbatim}
\newenvironment{reference}{\comment}{\endcomment}
%\newenvironment{reference}{}{}
\newenvironment{slogan}{\comment}{\endcomment}
\newenvironment{history}{\comment}{\endcomment}

% For commutative diagrams we use Xy-pic
\usepackage[all]{xy}

% We use 2cell for 2-commutative diagrams.
\xyoption{2cell}
% Theorem environments.
%
\theoremstyle{plain}
\newtheorem{theorem}[subsection]{Theorem}
\newtheorem{proposition}[subsection]{Proposition}
\newtheorem{lemma}[subsection]{Lemma}

\theoremstyle{definition}
\newtheorem{definition}[subsection]{Definition}
\newtheorem{example}[subsection]{Example}
\newtheorem{exercise}[subsection]{Exercise}
\newtheorem{situation}[subsection]{Situation}

\theoremstyle{remark}
\newtheorem{remark}[subsection]{Remark}
\newtheorem{remarks}[subsection]{Remarks}

\numberwithin{equation}{subsection}

% Macros
%
\def\lim{\mathop{\mathrm{lim}}\nolimits}
\def\colim{\mathop{\mathrm{colim}}\nolimits}
\def\Spec{\mathop{\mathrm{Spec}}}
\def\Hom{\mathop{\mathrm{Hom}}\nolimits}
\def\Ext{\mathop{\mathrm{Ext}}\nolimits}
\def\SheafHom{\mathop{\mathcal{H}\!\mathit{om}}\nolimits}
\def\SheafExt{\mathop{\mathcal{E}\!\mathit{xt}}\nolimits}
\def\Sch{\mathit{Sch}}
\def\Mor{\mathop{\mathrm{Mor}}\nolimits}
\def\Ob{\mathop{\mathrm{Ob}}\nolimits}
\def\Sh{\mathop{\mathit{Sh}}\nolimits}
\def\NL{\mathop{N\!L}\nolimits}
\def\CH{\mathop{\mathrm{CH}}\nolimits}
\def\proetale{{pro\text{-}\acute{e}tale}}
\def\etale{{\acute{e}tale}}
\def\QCoh{\mathit{QCoh}}
\def\Ker{\mathop{\mathrm{Ker}}}
\def\Im{\mathop{\mathrm{Im}}}
\def\Coker{\mathop{\mathrm{Coker}}}
\def\Coim{\mathop{\mathrm{Coim}}}

% Boxtimes
%
\DeclareMathSymbol{\boxtimes}{\mathbin}{AMSa}{"02}

%
% Macros for moduli stacks/spaces
%
\def\QCohstack{\mathcal{QC}\!\mathit{oh}}
\def\Cohstack{\mathcal{C}\!\mathit{oh}}
\def\Spacesstack{\mathcal{S}\!\mathit{paces}}
\def\Quotfunctor{\mathrm{Quot}}
\def\Hilbfunctor{\mathrm{Hilb}}
\def\Curvesstack{\mathcal{C}\!\mathit{urves}}
\def\Polarizedstack{\mathcal{P}\!\mathit{olarized}}
\def\Complexesstack{\mathcal{C}\!\mathit{omplexes}}
% \Pic is the operator that assigns to X its picard group, usage \Pic(X)
% \Picardstack_{X/B} denotes the Picard stack of X over B
% \Picardfunctor_{X/B} denotes the Picard functor of X over B
\def\Pic{\mathop{\mathrm{Pic}}\nolimits}
\def\Picardstack{\mathcal{P}\!\mathit{ic}}
\def\Picardfunctor{\mathrm{Pic}}
\def\Deformationcategory{\mathcal{D}\!\mathit{ef}}
\usepackage{cleveref}
\crefname{originalreference}{original reference}{original references}
\Crefname{originalreference}{Original reference}{Original references}
\makeatletter
\newlabel{examples-stacks-subsection-fppf-torsors}{{04UR}{source}{}{sourcepointer-1}{}}
\newlabel{examples-stacks-subsection-fppf-torsors@cref}{{[originalreference][0][]04UR}{[1][0][]source}{}{}{}}
\makeatother
\crefname{originalreference}{original reference}{original references}
\Crefname{originalreference}{Original reference}{Original references}
\makeatletter
\newlabel{examples-stacks-subsection-principal-homogeneous-spaces}{{04UN}{source}{}{sourcepointer-2}{}}
\newlabel{examples-stacks-subsection-principal-homogeneous-spaces@cref}{{[originalreference][0][]04UN}{[1][0][]source}{}{}{}}
\makeatother
\crefname{originalsection}{original section}{original sections}
\Crefname{originalsection}{Original section}{Original sections}
\makeatletter
\newlabel{groupoids-example-roots-of-unity}{{040M}{source}{}{sourcepointer-3}{}}
\newlabel{groupoids-example-roots-of-unity@cref}{{[originalsection][0][]040M}{[1][0][]source}{}{}{}}
\makeatother
\crefname{originalremark}{original remark}{original remarks}
\Crefname{originalremark}{Original remark}{Original remarks}
\makeatletter
\newlabel{groupoids-remark-fun-with-torsors}{{049C}{source}{}{sourcepointer-4}{}}
\newlabel{groupoids-remark-fun-with-torsors@cref}{{[originalremark][0][]049C}{[1][0][]source}{}{}{}}
\makeatother
\crefname{originallemma}{original lemma}{original lemmas}
\Crefname{originallemma}{Original lemma}{Original lemmas}
\makeatletter
\newlabel{varieties-lemma-locally-finite-type-Jacobson}{{0478}{source}{}{sourcepointer-5}{}}
\newlabel{varieties-lemma-locally-finite-type-Jacobson@cref}{{[originallemma][0][]0478}{[1][0][]source}{}{}{}}
\makeatother
\title{G-Principal and G-Torsors stacks are not equivalent in general}
\author{The Stacks Project Authors}
\date{stacks-project:examples.tex:local-snapshot-20261012; GNU FDL 1.2 or later}
\begin{document}
\maketitle
\noindent\textit{Editorial note.} All delivered mathematical text is the author's own, copied from the pinned Stacks Project chapter file examples.tex; no statement or proof step was rewritten, reordered or removed.\par\medskip
\noindent\textit{Editorial note.} The counted proof is the author's complete argument: The counted proof is the one-line conclusion of the section's explicit Q-example; the author's omitted classification proof is not load-bearing for this card, content verbatim.\par\medskip
\noindent\textit{Editorial note.} The label lemma-torsors-principal-spaces-not-equal names the counted statement, and the proof environment immediately following it in the source is delivered verbatim as its proof.\par\medskip
\noindent\textit{Editorial note.} The counted statement is carried by the author's own lemma environment in the chapter source.\par\medskip
\noindent\textit{Allowed standard prerequisites (editorial).}\begin{itemize}
\item Context span 4040-4096: G = (mu\_\{2,k\})\textasciicircum{}infty = lim (mu\_\{2,k\})\textasciicircum{}n, an affine group scheme with k[t\_1, t\_2, ...]/(t\_i\textasciicircum{}2 - 1); the torsor P = Spec(k[x\_1, x\_2, ...]/(x\_i\textasciicircum{}2 - a\_i)) with G-action x\_i -> t\_i  tensor  x\_i, which is a G-principal space (over k' = k(sqrt a\_1, ...) it becomes G-equivariantly isomorphic to G) and is an fppf torsor if and only if k'/k is finite (forward via the fppf covering Spec(k'), backward via a closed point of a trivilizing cover and Varieties, lemma-locally-finite-type-Jacobson); explicit example k = Q, a\_i = i gives k'/k infinite
\item Examples of Stacks subsections defining G-Principal and G-Torsors
\end{itemize}\medskip
\ifdefined\hypertarget\hypertarget{sourcepointer-1}{}\fi
\noindent\textit{Source reference.} Original reference 04UR: sources/stacks-chow/tags.txt, source lines 6299--6299. Prerequisite pointer invoked at 4105; it names a standard fact of a sibling chapter that the counted argument uses, and every step of the counted proof built on it is retained verbatim in the delivered body.\par\medskip
\ifdefined\hypertarget\hypertarget{sourcepointer-2}{}\fi
\noindent\textit{Source reference.} Original reference 04UN: sources/stacks-chow/tags.txt, source lines 6296--6296. Prerequisite pointer invoked at 4102; it names a standard fact of a sibling chapter that the counted argument uses, and every step of the counted proof built on it is retained verbatim in the delivered body.\par\medskip
\ifdefined\hypertarget\hypertarget{sourcepointer-3}{}\fi
\noindent\textit{Source reference.} Original section 040M: sources/stacks-chow/tags.txt, source lines 5261--5261. Prerequisite pointer invoked at 4130; it names a standard fact of a sibling chapter that the counted argument uses, and every step of the counted proof built on it is retained verbatim in the delivered body.\par\medskip
\ifdefined\hypertarget\hypertarget{sourcepointer-4}{}\fi
\noindent\textit{Source reference.} Original remark 049C: sources/stacks-chow/tags.txt, source lines 5575--5575. Prerequisite pointer invoked at 4116; it names a standard fact of a sibling chapter that the counted argument uses, and every step of the counted proof built on it is retained verbatim in the delivered body.\par\medskip
\ifdefined\hypertarget\hypertarget{sourcepointer-5}{}\fi
\noindent\textit{Source reference.} Original lemma 0478: sources/stacks-chow/tags.txt, source lines 5499--5499. Prerequisite pointer invoked at 4162; it names a standard fact of a sibling chapter that the counted argument uses, and every step of the counted proof built on it is retained verbatim in the delivered body.\par\medskip
\noindent
In
Groupoids, Remark \ref{groupoids-remark-fun-with-torsors}
we raise the question whether any $G$-torsor is a $G$-torsor for the
fppf topology. In this section we show that this is not always the case.

\medskip\noindent
Let $k$ be a field. All schemes and stacks are over $k$ in what follows.
Let $G \to \Spec(k)$ be the group scheme
$$
G = (\mu_{2, k})^\infty =
\mu_{2, k} \times_k \mu_{2, k} \times_k \mu_{2, k} \times_k \ldots =
\lim_n (\mu_{2, k})^n
$$
where $\mu_{2, k}$ is the group scheme of second roots of unity over
$\Spec(k)$, see
Groupoids, Example \ref{groupoids-example-roots-of-unity}.
As an inverse limit of affine schemes we see that $G$ is an affine group
scheme. In fact it is the spectrum of the ring
$k[t_1, t_2, t_3, \ldots]/(t_i^2 - 1)$. The multiplication map
$m : G \times_k G \to G$ is on the algebra level given by
$t_i \mapsto t_i \otimes t_i$.

\medskip\noindent
We claim that any $G$-torsor over $k$ is of the form
$$
P = \Spec(k[x_1, x_2, x_3, \ldots]/(x_i^2 - a_i))
$$
for certain $a_i \in k^*$ and with $G$-action $G \times_k P \to P$
given by $x_i \to t_i \otimes x_i$ on the algebra level.
We omit the proof.
Actually for the example we only need that $P$ is a $G$-torsor
which is clear since over $k' = k(\sqrt{a_1}, \sqrt{a_2}, \ldots)$
the scheme $P$ becomes isomorphic to $G$ in a $G$-equivariant manner.
Note that $P$ is trivial if and only if $k' = k$ since if
$P$ has a $k$-rational point then all of the $a_i$ are squares.

\medskip\noindent
We claim that $P$ is an fppf torsor if and only if the field extension
$k' = k(\sqrt{a_1}, \sqrt{a_2}, \ldots)/k$ is finite.
If $k'$ is finite over $k$, then
$\{\Spec(k') \to \Spec(k)\}$
is an fppf covering which trivializes $P$ and we see that $P$ is indeed
an fppf torsor. Conversely, suppose that $P$ is an fppf $G$-torsor.
This means that there exists an fppf covering
$\{S_i \to \Spec(k)\}$ such that each $P_{S_i}$ is trivial.
Pick an $i$ such that $S_i$ is not empty. Let $s \in S_i$ be a closed
point. By
Varieties, Lemma \ref{varieties-lemma-locally-finite-type-Jacobson}
the field extension $\kappa(s)/k$ is finite, and by construction
$P_{\kappa(s)}$ has a $\kappa(s)$-rational point. Thus we see that
$k \subset k' \subset \kappa(s)$ and $k'$ is finite over $k$.

\medskip\noindent
To get an explicit example take $k = \mathbf{Q}$ and $a_i = i$
for example (or $a_i$ is the $i$th prime if you like).


\begin{lemma}
\label{lemma-torsors-principal-spaces-not-equal}
Let $S$ be a scheme. Let $G$ be a group scheme over $S$.
The stack $G\textit{-Principal}$ classifying principal homogeneous $G$-spaces
(see Examples of Stacks, Subsection
\ref{examples-stacks-subsection-principal-homogeneous-spaces})
and the stack $G\textit{-Torsors}$ classifying fppf $G$-torsors
(see Examples of Stacks, Subsection
\ref{examples-stacks-subsection-fppf-torsors})
are not equivalent in general.
\end{lemma}

\begin{proof}
The discussion above shows that the functor
$G\textit{-Torsors} \to G\textit{-Principal}$ isn't essentially
surjective in general.
\end{proof}

\end{document}
